\documentclass[10pt,a4paper]{amsart}
\usepackage[margin=19mm]{geometry}
\usepackage{amsmath,amssymb,mathtools}
\usepackage[scr=rsfs,cal=euler]{mathalfa}
\usepackage{xfrac}
\usepackage[unicode,hidelinks,bookmarks=false]{hyperref}
\newcommand{\ie}{i.e.\ }
\newcommand{\cf}{cf.\ }
\newcommand{\resp}{respectively\ }
\newcommand{\Subsection}{\subsection}
\providecommand{\nobreakdash}{\nobreak\hbox{-}}
\setlength{\emergencystretch}{2em}
\multlinegap0pt
\usepackage{amssymb}
\usepackage{cancel}
\usepackage{enumerate}
\usepackage[shortlabels]{enumitem}
\usepackage{bm}
\usepackage{tikz}
\newtheorem{proposition}{Proposition}
\newtheorem{theorem}{Theorem}
\title{Various Properties of Various Ultrafilters, Various Graph Width Parameters, and Various Connectivity Systems (with Survey)}
\author{Takaaki Fujita}
\date{Source: arXiv:2408.02299v9 (CC BY 4.0)}
\begin{document}
\maketitle

\noindent\emph{Source and licence.} Various Properties of Various Ultrafilters, Various Graph Width Parameters, and Various Connectivity Systems (with Survey). Version arXiv:2408.02299v9 (CC BY 4.0), distributed by arXiv under the Creative Commons Attribution 4.0 International licence, http://creativecommons.org/licenses/by/4.0/, as recorded in the arXiv licence metadata for this exact version and shown on the abstract page https://arxiv.org/abs/2408.02299v9. Full licence notice: \texttt{licenses/arxiv-2408.02299-CC-BY-4.0.txt}.\par\smallskip

\noindent\textit{Editorial scope.} Counted unit: the article's theorem thm:filter\_generated\_by\_subbase (source0.tex lines 2243--2247). The article's complete original proof of it is delivered with it (source0.tex lines 2249--2281, one \texttt{proof} environment of 33 lines). No mathematics is written, repaired or re-derived: the statement and the proof are the article's own lines, in the article's own notation, and no retained line is substituted or re-flowed.  Retained uncounted as the source-local context the proof needs: lines 2196--2200.  Nothing was omitted from any retained span, so the package carries no omission record.  No retained line contains a citation, so the wrapper carries no bibliography and no bibliography entry.  No retained line contains a figure environment, an image-inclusion command or drawing code, so no image is delivered and the package declares no asset dependency.  Editorial formatting only, in addition: an AMS \texttt{amsart} preamble that restates the article's own declarations and macros used by the retained lines, namely the environments \texttt{proposition}, \texttt{theorem} on separate counters, together with the standard packages those lines need.  The retained environments print in this document as Proposition 1, Theorem 1 in order of appearance, whereas the article prints the same items with the numbering of its own full document.  The complete native source of the article as distributed in the arXiv e-print for this exact version is bundled unchanged as \texttt{sources/163346/source0.tex}.
\begin{proposition}
\label{thm:prefilter_generated_by_subbase}
Let \(S\) be a filter subbase of order \(k+1\) on a connectivity system \((X,f)\).
Then \(\operatorname{Pf}(S)\) is a prefilter of order \(k+1\).
\end{proposition}

\begin{theorem}
\label{thm:filter_generated_by_subbase}
Let \(S\) be a filter subbase of order \(k+1\) on a connectivity system \((X,f)\).
Then \(\operatorname{Fil}(S)\) is a filter of order \(k+1\).
\end{theorem}

\begin{proof}
We verify axioms \((Q0)\)--\((Q3)\).

\emph{(Q0)} holds by definition of \(\operatorname{Fil}(S)\).

To show nonemptiness, choose \(A\in \operatorname{Pf}(S)\).
Then \(A\subseteq A\) and \(f(A)\leq k\), so \(A\in \operatorname{Fil}(S)\).
Thus \(\operatorname{Fil}(S)\neq \emptyset\).

\emph{(Q3)}:
Since every \(A\in \operatorname{Pf}(S)\) is nonempty by (SB4), the empty set cannot contain any member of \(\operatorname{Pf}(S)\).
Hence \(\emptyset\notin \operatorname{Fil}(S)\).

\emph{(Q2)}:
Let \(B\in \operatorname{Fil}(S)\), and suppose \(B\subseteq C\subseteq X\) with \(f(C)\leq k\).
Choose \(A\in \operatorname{Pf}(S)\) such that \(A\subseteq B\).
Then \(A\subseteq C\), so \(C\in \operatorname{Fil}(S)\).

\emph{(Q1)}:
Let \(B,C\in \operatorname{Fil}(S)\), and assume \(f(B\cap C)\leq k\).
Choose \(A_B,A_C\in \operatorname{Pf}(S)\) such that
\[
A_B\subseteq B,\qquad A_C\subseteq C.
\]
Since \(\operatorname{Pf}(S)\) is a prefilter by Proposition \ref{thm:prefilter_generated_by_subbase},
there exists \(A\in \operatorname{Pf}(S)\) such that
\[
A\subseteq A_B\cap A_C \subseteq B\cap C.
\]
Because \(f(B\cap C)\leq k\), it follows that \(B\cap C\in \operatorname{Fil}(S)\).

Therefore \(\operatorname{Fil}(S)\) is a filter of order \(k+1\).
\end{proof}

\end{document}
